\usepackage{etoolbox} % toggle
\usepackage[capitalize,noabbrev,sort]{cleveref} % cref
\usepackage{listings}
\usepackage{subfig}
\usepackage{wrapfig}
\usepackage{listings}
\usepackage{url}
\usepackage{algorithm}
\usepackage{textcomp}
\usepackage{xspace}
\usepackage{hyperref}
\usepackage{mathpartir}
\usepackage{natbib}
\usepackage{diagbox}
\usepackage[normalem]{ulem} % needed for sout

\usepackage{tikz}
\usetikzlibrary{tikzmark,shapes.misc,decorations.pathmorphing}
\usepackage{mathtools}

%\usepackage[dvipsnames]{xcolor}
%\usepackage{lstparams, lstcoq}

\usepackage[export]{adjustbox} % makebox
\usepackage{multirow}
\usepackage{hhline}
\usepackage{pifont} % ding
\usepackage[inline]{enumitem}
\setlist{leftmargin=5.5mm,topsep=2mm}
\setlist[enumerate]{label=\textbf{(\arabic*)}}

\sloppypar

%%%%%%%%%%%%%%%%%%%%%%%%%%%%%%%%%%%%%%%%%%%%%%%%%%%%%%%%%%%%%%%%%%%%%%%%%%%%%%%%
% Comments
%%%%%%%%%%%%%%%%%%%%%%%%%%%%%%%%%%%%%%%%%%%%%%%%%%%%%%%%%%%%%%%%%%%%%%%%%%%%%%%%

\newtoggle{comments}
\toggletrue{comments}
%\togglefalse{comments}

\setlength{\marginparwidth}{47pt}
\setlength{\marginparsep}{3pt}
\newcommand{\mpfont}{\scriptsize\sffamily}
\newcommand{\MPworker}[2]{
  \unskip{\color{#1}\vrule\vrule}
  {\marginpar{\raggedright\color{#1}\mpfont #2}}
}
\newcommand{\todo}[1]{
  \iftoggle{comments} {\MPworker{blue}{To-do: #1}} {}
}

%\newcommand{\todo}[1]{
%    \iftoggle{comments} {{\color{teal}[TODO: #1]}} {}
%}
\newcommand{\cancut}[1]{
  %    \iftoggle{comments} {{\color{orange}#1}} {#1}
  %{{\color{orange}#1}}
  {}
}

\newcommand{\namecomment}[3]{
  \iftoggle{comments} {\textbf{\textcolor{#2}{[#1: #3]}}} {}
}

% ADD YOUR NAME FOLLOWING THIS FORMAT
% Default color set:
%  red, green, blue, cyan, magenta, yellow, black, gray, darkgray, lightgray,
%  brown, lime, olive, white, orange, pink, purple, teal, violet.
\newcommand{\jiyong}[1]{\namecomment{Ji-Yong}{cyan}{#1}}
\newcommand{\toyifei}[1]{\namecomment{ATTN, YIFEI:}{red}{#1}}
\newcommand{\yifei}[1]{\namecomment{Yifei}{teal}{#1}}

%%%%%%%%%%%%%%%%%%%%%%%%%%%%%%%%%%%%%%%%%%%%%%%%%%%%%%%%%%%%%%%%%%%%%%%%%%%%%%%%
% Place holders
%%%%%%%%%%%%%%%%%%%%%%%%%%%%%%%%%%%%%%%%%%%%%%%%%%%%%%%%%%%%%%%%%%%%%%%%%%%%%%%%

\newcommand{\needcite}[1]{\textbf{\textcolor{red}{[cite: #1]}}}
\newcommand{\needref}[1]{\textbf{\textcolor{red}{[ref: #1]}}}
\newcommand{\ignore}[1]{}

%%%%%%%%%%%%%%%%%%%%%%%%%%%%%%%%%%%%%%%%%%%%%%%%%%%%%%%%%%%%%%%%%%%%%%%%%%%%%%%%
% Custom names
%%%%%%%%%%%%%%%%%%%%%%%%%%%%%%%%%%%%%%%%%%%%%%%%%%%%%%%%%%%%%%%%%%%%%%%%%%%%%%%%
\newcommand{\sysname}{\textsc{SemantiX}\xspace}
\newcommand{\sysnamelong}{semantic compatibility checking framework
between applications and composition of distributed systems}

\newcommand{\sysgraph}{AppGraph\xspace}
\newcommand{\sysgraphs}{AppGraphs\xspace}

\newcommand{\sysmod}{DM\xspace}
\newcommand{\sysmods}{DMs\xspace}
\newcommand{\sysmodlong}{distributed system module}
\newcommand{\sysmodlongs}{distributed system modules}

\newcommand{\eg}{\emph{e.g.}, }
\newcommand{\ie}{\emph{i.e.}, }

\definecolor{StringRed}{rgb}{.637,0.082,0.082}
\definecolor{CommentGreen}{rgb}{0.0,0.55,0.3}
\newcommand{\greenCheck}{\textcolor{CommentGreen}{\ding{51}}}%
\newcommand{\redX}{\textcolor{StringRed}{\ding{55}}}%
\newcommand{\yellowTriangle}{\textcolor{olive}{\ding{115}}}%
\newcommand{\purpleTriangle}{\textcolor{purple}{\ding{115}}}%

\newcommand{\wcmark}{{\color{black}\ding{52}}}
\newcommand{\cmark}{{\color{blue}\ding{52}}}
\newcommand{\xmark}{{\color{red}\ding{56}}}

%%%%%%%%%%%%%%%%%%%%%%%%%%%%%%%%%%%%%%%%%%%%%%%%%%%%%%%%%%%%%%%%%%%%%%%%%%%%%%%%
% Custom enumerate/itemize/dictionary formats
%%%%%%%%%%%%%%%%%%%%%%%%%%%%%%%%%%%%%%%%%%%%%%%%%%%%%%%%%%%%%%%%%%%%%%%%%%%%%%%%
\usepackage{enumitem}

\newcommand{\enumparen}[1]{
  \begin{enumerate}[noitemsep,topsep=0pt]
      {#1}
  \end{enumerate}
}
\newcommand{\itemparen}[1]{
  \begin{itemize}[noitemsep,topsep=0pt]
      {#1}
  \end{itemize}
}
\newcommand{\customenumparen}[3]{
  \begin{enumerate}[label={#1},noitemsep,topsep=0pt,left=5pt..25pt]
      \setcounter{enumi}{#2}
      {#3}
  \end{enumerate}
}

%%%%%%%%%%%%%%%%%%%%%%%%%%%%%%%%%%%%%%%%%%%%%%%%%%%%%%%%%%%%%%%%%%%%%%%%%%%%%%%%
% Other formats
%%%%%%%%%%%%%%%%%%%%%%%%%%%%%%%%%%%%%%%%%%%%%%%%%%%%%%%%%%%%%%%%%%%%%%%%%%%%%%%%
\newcommand{\ballnumber}[1]{
  \tikz[baseline=(myanchor.base)]
  \node[circle,fill=.,inner sep=1pt] (myanchor)
  {\color{-.}\bfseries\footnotesize #1};
}

\newcommand{\circled}[1]{
  \tikz[baseline=(char.base)]{
    \node[shape=circle,draw,inner sep=1pt] (char)
    {#1};
  }
}

\newcommand{\para}[1]{\paragraph{#1}}
\newcommand{\subpara}[1]{\subparagraph{#1}}

%\newtheorem{assumption}{Assumption}

%%%%%%%%%%%%%%%%%%%%%%%%%%%%%%%%%%%%%%%%%%%%%%%%%%%%%%%%%%%%%%%%%%%%%%%%%%%%%%%%
% Fine grain edits
%%%%%%%%%%%%%%%%%%%%%%%%%%%%%%%%%%%%%%%%%%%%%%%%%%%%%%%%%%%%%%%%%%%%%%%%%%%%%%%%

% I'd rather add package "changes" but it conflicts with other
% things, so this is good enough.
% Prefix "sug" means "suggested"
\definecolor{SuggestionGreen}{HTML}{35713A} % taken from Google Docs
\newcommand{\suginsert}[1]{\textbf{\textcolor{SuggestionGreen}{#1}}}
% see https://tex.stackexchange.com/a/130626
\newcommand\sugdelete{\bgroup\markoverwith{\textcolor{SuggestionGreen}{\rule[0.35ex]{2pt}{0.35ex}}}\ULon}
\newcommand{\sugreplace}[2]{\suginsert{#1} \sugdelete{#2}}

%%%%%%%%%%%%%%%%%%%%%%%%%%%%%%%%%%%%%%%%%%%%%%%%%%%%%%%%%%%%%%%%%%%%%%%%%%%%%%%%
% Custom section/paragraph formats
%%%%%%%%%%%%%%%%%%%%%%%%%%%%%%%%%%%%%%%%%%%%%%%%%%%%%%%%%%%%%%%%%%%%%%%%%%%%%%%%
\usepackage{titlesec}

%-------------------------------------------------------------------------------
% \titleformat{command}[shape]{format}{label}{space}{before code}[after code]
%-------------------------------------------------------------------------------
% command: target command
% shape: [runin: embed title in the paragraph] [hang: default as in sections]
% format: title text format
% labeling: numbering format for table of contents
% space: space after title
% before code: code betore title
% after code: code after title

%\titleformat{\paragraph}[runin]{\normalfont\normalsize\bfseries}{\theparagraph}{0em}{\uline}
\titleformat{\paragraph}[runin]{\normalfont\normalsize\bfseries\itshape}{\theparagraph}{0em}{}

%\titleformat{\subparagraph}[runin]{\normalfont\normalsize\bfseries\itshape}{\thesubparagraph}{1em}{}
%\titleformat{\subparagraph}[runin]{\normalfont\normalsize}{\thesubparagraph}{1em}{\underline}

%-------------------------------------------------------------------------------
% \titlespacing{command}{leftmargin}{before space}{after space}[right space]
%-------------------------------------------------------------------------------
\titlespacing*{\paragraph}{0pt}{5pt}{0.5em}
%\titlespacing*{\subparagraph}{\parindent}{5pt}{1em}

\newcommand*{\squigglyarrowhead}[0]{\ \raisebox{0.15em}{\tikz \draw [->,
      line join=round,
      line width=0.8pt,
      decorate, decoration={
        zigzag,
        segment length=6,
        amplitude=.9,
        post=lineto,
        post length=1pt
}] (0,0) -- (0.4,0);}\ }
\newcommand*{\squigglyarrowtail}[0]{\ \raisebox{0.15em}{\tikz \draw [-,
      line join=round,
      line width=0.8pt,
      decorate, decoration={
        zigzag,
        segment length=6,
        amplitude=.9,
        post=lineto,
        post length=1pt
}] (0,0) -- (0.4,0);}\ }

\newcommand{\ltrans}[1]{\mathopen{\squigglyarrowtail} \hspace{-3pt}
\fbox{$#1$} \hspace{-3pt} \mathclose{\squigglyarrowhead}}
\newcommand{\lstep}[3]{#1 \mathop{\ltrans{#2}} #3}
\newcommand{\lstepStar}[3]{#1 \mathop{\ltrans{#2} \hspace{-3pt}^*} #3}
\newcommand{\llabel}[1]{\fbox{$#1$}}

\DeclareMathOperator*{\eventually}{\Diamond}
\DeclareMathOperator*{\always}{\Box}
